\chapter{Коэрцитивность полного высокоуглового тензорного хвоста}
\label{ch:v6-bosonic-defect-full-tensor-high-angular-coercivity-gate}

\section{Постановка}

Переносная калибровка сняла ложные отрицательные пары и проверила только
окно $-3\le j\le3$. Для сертификата прямой нити требуется покрыть все
целые угловые метки каждого из трёх скрученных характеров $Z_3$. Как и в
стандартном частично-волновом анализе вихрей, доказательство разделяется на
конечный численный мост и центробежно коэрцитивный хвост
\cite{full-tensor-high-angular-baacke}.

\section{Повторная проверка низкого окна}

Чтобы исключить пропуск отрицательного уровня при поиске около нуля, блоки
$c=-1,0,1$, $-3\le j\le3$ повторно решены плотным обобщённым эрмитовым
алгоритмом, возвращающим алгебраически минимальное собственное значение.
Отрицательных уровней нет. Переносная пара при $(c,j)=(0,\pm1)$ равна
$0.00145554$ на сетке $80$ и стремится к нулю при сгущении.

После удаления этой пары нижний внутренний уровень равен
\begin{equation}
 \lambda_{\mathrm{int},80}=3.61877141
\end{equation}
и лежит в сопряжённых блоках $(c,j)=(1,1),(-1,-1)$. Точное континуальное
значение этой щели будет предметом отдельного гейта.

\section{Прямой мост}

На той же сетке плотным решателем проверены все знаки, три характера и
метки
\begin{equation}
 4\le |j|\le139.
\end{equation}
Глобальный минимум моста достигается при $|j|=4$:
\begin{equation}
 \lambda_{4}=3.77339594.
\end{equation}
Минимум по характерам и знакам строго возрастает с $|j|$; характерные
значения равны
\begin{equation}
 \lambda_{10}=4.14123894,
 \quad
 \lambda_{50}=11.66223455,
 \quad
 \lambda_{100}=32.66358434,
 \quad
 \lambda_{139}=57.41164870.
\end{equation}
Отрицательных алгебраических уровней в мосте нет, максимальный эрмитов
остаток равен $3.18\cdot10^{-14}$.

\section{Аналитический хвост}

Для каждого характера и обоих знаков $j$ построена метрически нормированная
локальная матрица размера $6\times6$. В ней сохранены:
\begin{enumerate}
 \item полный угловой квадрат четырёх тензорных компонент;
 \item половина полярного квадрата Ходжа связности;
 \item потенциальный гессиан полного поля $Q+T$;
 \item радиальные и угловые перекрёстные члены второй вариации;
 \item неабелев перекрёстный член кривизны;
 \item штраф источника фоновой калибровки после неравенства Юнга.
\end{enumerate}

Каждая строковая оценка Гершгорина имеет вид
\begin{equation}
 g_{c,\sigma,r,k}(J)=aJ^2+bJ+d,
 \qquad J=|j|.
\end{equation}
На $1599$ радиальных элементах найдено
\begin{equation}
 \min a=0.0012510559>0.
\end{equation}
Первая положительная выборочная граница возникает при $J=130$. Для
консервативного сертификата выбран порог $J=140$, где
\begin{equation}
 \min g(140)=1.21178043,
 \qquad
 \min g'(140)=0.12813910.
\end{equation}
Следовательно, нижняя оценка положительна и возрастает для всех целых
$|j|\ge140$.

\begin{theorem}
Калиброванный полный полярный гессиан стационарного прямого вихря не имеет
отрицательных мод ни в одном из трёх скрученных характеров и ни при одной
целой угловой метке. После выделения переносной пары оператор линейно
положителен; точное континуальное значение внутренней щели ещё не вычислено.
\end{theorem}

\begin{nogocriterion}
Теорема относится к бесконечной прямой нити и поперечному оператору.
Продольные волны добавляют неотрицательный квадрат импульса, но замыкание
нити в петлю вводит кривизну и глобальные моды. Поэтому устойчивость петли и
рождение локализованной частицы из результата не следуют автоматически.
\end{nogocriterion}

\section{Следующий различающий тест}

Нулевая гипотеза: сопряжённая внутренняя пара $(1,1),(-1,-1)$ сходится к
строго положительной континуальной щели и остаётся ниже остальных внутренних
каналов.

Стоп-критерий: если щель уходит к нулю после удаления переносов, прямая нить
имеет дополнительный модуль и не является изолированным устойчивым объектом.
Если предел положителен, поперечная линейная устойчивость прямой нити будет
закрыта вместе с численным значением щели.

\begin{protocol}
\begin{enumerate}
 \item плотная проверка $|j|\le3$: \textbf{пройдена};
 \item прямой мост $4\le|j|\le139$: \textbf{пройден};
 \item характеры $c=-1,0,1$: \textbf{покрыты};
 \item оба знака $j$: \textbf{покрыты};
 \item аналитический порог: \textbf{$|j|=140$};
 \item отрицательные внутренние моды: \textbf{не найдены};
 \item полный высокоугловой хвост: \textbf{закрыт};
 \item точная внутренняя щель: \textbf{открыта};
 \item замкнутая петля: \textbf{не проверена};
 \item рождение материи: \textbf{не закрыто}.
\end{enumerate}
\end{protocol}

Машинный сертификат сохранён в
\path{s2t_v6_bosonic_defect_full_tensor_high_angular_coercivity_gate_results.json}.

\begin{thebibliography}{9}
\bibitem{full-tensor-high-angular-baacke}
J.~Baacke, N.~Kevlishvili,
\emph{One-loop corrections to the string tension of the vortex in the
Abelian Higgs model}, arXiv:0806.4349.
\end{thebibliography}